\chapter{Binding Notions and Their Bounds}
\label{ch:corollaries}

This chapter derives, from \cref{thm:symbolic-list,thm:identification}, the bounds for
six security notions of the symbolic protocol: unique binding, proximity, promise
binding, session uniqueness, anchored binding and non-equivocation. Each notion is stated as a game on
the symbolic prover of \cref{def:symbolic-game}. Throughout,
\[
\theta=\tfrac{1-\rho}2,\qquad
\eps_{\mathrm{UD}}:=\frac{N-B}{|\K|},\qquad
\eps_{\mathrm{SZ}}(\ell):=\Bigl(\ell+\tbinom{\ell}2\Bigr)\frac{d}{|\K|-1}.
\]

%=====================================================================
\section{Unique binding}
%=====================================================================

\begin{definition}[Unique binding game]\label{def:game-A}
Let $F_\star:=F_\star(W_0)$ be the reference polynomial of \cref{def:reference} and,
when $F_\star\ne\bot$, $f_\star:=f_{F_\star}=\kappa^{-1}(F_\star)$. The prover wins if
\[
\Acc\ \wedge\ \bigl(F_\star=\bot\ \vee\ \gt(T)\ne f_\star(r_1T,\dots,r_kT)\bigr).
\]
\end{definition}

By \cref{def:reference}, $F_\star\in\F[X]$ and hence $f_\star\in\Pa(\F)$: the reference
polynomial has coefficients in the base field.

\begin{theorem}[Unique binding]\label{thm:unique-binding}
For every symbolic prover and every arity sequence,
\[
\Pr[\text{the prover wins the unique binding game}]
\le\frac{N-B}{|\K|}+\frac{d}{|\K|-1}+\Bigl(\frac{1+\rho}2\Bigr)^{s}.
\]
\end{theorem}

\begin{proof}
Let $\alpha$ be the smallest element of $\frac1{n_1}\Z$ strictly greater than
$\frac{1+\rho}2$ and $\delta:=1-\alpha<\theta$. Then $\alpha^-\le\frac{1+\rho}2$. If
$\agr(W_0,G)\ge\alpha$ then $\Db(W_0,G)<\theta$, so
$\Lambda_\alpha(W_0)\subseteq\{F_\star\}\setminus\{\bot\}$ and $\ell_\alpha\le1$. The
winning event is therefore contained in the event of \cref{thm:symbolic-list}, which is
bounded by $\eps_{\MCA}(\delta)+d/(|\K|-1)+(\alpha^-)^s$. Finally
$\eps_{\MCA}(\delta)\le\eps_{\mathrm{UD}}$ by \cref{cor:mca-ud-sum}.
\end{proof}

\begin{remark}\label{rem:unique-binding-structure}
The bound depends on the arities only through $d=\sum_i(a_i-1)$, which enters a term of
order $d/|\K|$. The query term depends on $\rho$ alone. The next chapter shows that
the factor $\frac{1+\rho}2$ cannot be improved for this game.
\end{remark}

%=====================================================================
\section{Proximity}
%=====================================================================

\begin{theorem}[Proximity]\label{thm:proximity}
Let $\alpha\in[\rho,1]$ and $\delta:=1-\alpha$. For every symbolic prover,
\[
\Pr\bigl[\Acc\ \wedge\ \Lambda_\alpha(W_0)=\emptyset\bigr]\le\eps_{\MCA}(\delta)+(\alpha^-)^s .
\]
In particular, for $\alpha=\sqrt\rho+\eta$ with $\eta>0$, the right-hand side is at most
\[
\sum_{i<k}\eps^J_{\MCA}(n_{i+1},\rho,a_{i+1},\eta,|\K|)+(\sqrt\rho+\eta)^s,
\]
and for
$\alpha$ the smallest element of $\frac1{n_1}\Z$ above $\frac{1+\rho}2$ it is at most
$\eps_{\mathrm{UD}}+(\frac{1+\rho}2)^s$.
\end{theorem}

\begin{proof}
If $\Lambda_\alpha(W_0)=\emptyset$ then the set $Z(\tau)$ in the proof of
\cref{thm:symbolic-list} is empty, so that proof gives the bound without the root-count
term. The two instances follow from \cref{thm:mca-johnson} and \cref{cor:mca-ud-sum}.
\end{proof}

%=====================================================================
\section{Promise binding}
%=====================================================================

\begin{lemma}[Lists under a proximity promise]\label{lem:promise-list}
Let $F\in\K[X]_{<L}$ with $\Db(W_0,F)\le\delta_0$, and let $\alpha>\delta_0+\rho$. Then
$\Lambda_\alpha(W_0)\subseteq\{F\}$.
\end{lemma}

\begin{proof}
Let $G\in\Lambda_\alpha(W_0)$. The fibres over which $W_0=G$ number at least $\alpha n_1$,
those over which $W_0=F$ at least $(1-\delta_0)n_1$; so $G=F$ on at least
$(\alpha-\delta_0)n_1>\rho n_1$ fibres. By \cref{lem:codeword-fibres}, $G=F$.
\end{proof}

\begin{theorem}[Promise binding]\label{thm:promise}
Let $\delta_0\in[0,1-\rho)$, let $\alpha\in(\delta_0+\rho,1]$ and $\delta:=1-\alpha$. For
every symbolic prover and every $F\in\K[X]_{<L}$,
\[
\Pr\bigl[\Acc\ \wedge\ \Db(W_0,F)\le\delta_0\ \wedge\ \gt\ne f_F(rT)\bigr]
\le\eps_{\MCA}(\delta)+\frac{d}{|\K|-1}+(\alpha^-)^s .
\]
\end{theorem}

\begin{proof}
On $\Db(W_0,F)\le\delta_0$, \cref{lem:promise-list} gives $\Lambda_\alpha(W_0)\subseteq\{F\}$,
so $\gt\ne f_F(rT)$ implies $\gt\notin\{f_G(rT):G\in\Lambda_\alpha(W_0)\}$. In the proof of
\cref{thm:symbolic-list} the list has at most one element, so the root-count term is
$d/(|\K|-1)$.
\end{proof}

\begin{remark}\label{rem:promise-johnson}
With MCA up to the Johnson radius, \cref{thm:promise} applies to every
$\alpha\ge\max(\sqrt\rho+\eta,\ \delta_0+\rho+\tfrac1{n_1})$. For $\delta_0<\sqrt\rho-\rho$ the
query factor is therefore $\sqrt\rho+\eta$ rather than $\frac{1+\rho}2$. The promise is a
property of the committed word and is not checked by the protocol.
\end{remark}

%=====================================================================
\section{Session uniqueness}
%=====================================================================

\begin{theorem}[Session uniqueness]\label{thm:session}
Let $\alpha\in[\rho,1]$ and $\delta:=1-\alpha$. For every symbolic prover,
\[
\Pr\bigl[\Acc\ \wedge\ G(\tau)=\bot\bigr]\le\eps_{\MCA}(\delta)+\eps_{\mathrm{SZ}}(\ell_\alpha)+(\alpha^-)^s,
\]
with $G(\tau)$ as in \cref{def:opened}. Outside this event, the accepted $\gt$ equals
$f_{G(\tau)}(rT)$ for a single $G(\tau)\in\Lambda_\alpha(W_0)$, and every value
$\gt(c'')$, $c''\in\K$, equals $f_{G(\tau)}(c''r)$.
\end{theorem}

\begin{proof}
If $\pi^{\mathrm{sym}}_0<\alpha$ then $\Pr[\Acc\mid\tau]\le(\alpha^-)^s$. If
$\pi^{\mathrm{sym}}_0\ge\alpha$ and $G(\tau)=\bot$, the event $\mathsf E$ of
\cref{thm:identification} occurs. The last sentence is \cref{rem:reuse}.
\end{proof}

%=====================================================================
\section{Anchored binding up to the Johnson radius}
%=====================================================================

Unique binding cannot go below the factor $\frac{1+\rho}2$ (\cref{thm:ub-tight}): beyond the
unique-decoding radius the committed word may be close to several polynomials, and the
prover may open any of them. An out-of-domain value fixed at commitment time selects one
list element. In the symbolic protocol this costs no extra query: the folding line is
chosen to pass through the Kronecker point of the out-of-domain sample, so the value is
read off as $\gt(0)$.

\begin{definition}[Kronecker point]\label{def:kronecker-point}
For $z\in\K$ let $\zeta(z):=\bigl(\Sigma_0(z),\Sigma_1(z),\dots,\Sigma_{k-1}(z)\bigr)\in\K^k$.
By \cref{def:kronecker}, $\kappa(f)(z)=f(\zeta(z))$ for every $f\in\Pa(\K)$, i.e.\
$G(z)=f_G(\zeta(z))$ for every $G\in\K[X]_{<L}$.
\end{definition}

\begin{protocol}[Anchored symbolic FRI]\label{prot:anchored}
\emph{Commitment.} $\mathsf P$ sends $W_0:D_0\to\F$; $\mathsf V$ samples $z\leftarrow\K$
uniformly; $\mathsf P$ sends $y\in\K$ (honestly $y=F(z)$ for $F=\kappa(f)$). The commitment is
$(W_0,z,y)$; put $\zeta:=\zeta(z)$.

\emph{Session.} Run \cref{prot:sfri} on $W_0$ with the challenges $r_{i}T$ replaced by
$\zeta_i+r_iT$ (\cref{def:curve-challenges} with $e=1$ and base point $\zeta$), and in
addition let $\mathsf V$ reject unless $\gt(0)=y$. Any number of sessions may be run on
the same commitment, each with fresh challenges and queries.
\end{protocol}

For an honest prover, \cref{prop:honest} with the curve challenges gives
$\gt=f(\zeta+rT)$, hence $\gt(0)=f(\zeta)=F(z)=y$: the protocol is complete.

\begin{lemma}[Out-of-domain separation]\label{lem:ood}
Let $W_0$ be fixed and $z\leftarrow\K$ uniform. The probability that two distinct elements of
$\Lambda_\alpha(W_0)$ take the same value at $z$ is at most
$\eps_{\mathrm{OOD}}(\ell_\alpha):=\binom{\ell_\alpha}2\frac{L-1}{|\K|}$.
\end{lemma}

\begin{proof}
For $G\ne G'$ in the list, $G-G'$ is nonzero of degree $<L$; union bound over pairs.
\end{proof}

\begin{definition}[Anchored reference]\label{def:anchored-ref}
For a commitment $(W_0,z,y)$ let $G_{z,y}$ be the unique $G\in\Lambda_\alpha(W_0)$ with
$G(z)=y$ if exactly one exists, and $G_{z,y}:=\bot$ otherwise. The prover wins a session if
\[
\Acc\ \wedge\ \bigl(G_{z,y}=\bot\ \vee\ \gt\ne f_{G_{z,y}}(\zeta+rT)\bigr).
\]
\end{definition}

\begin{theorem}[Anchored binding]\label{thm:anchored}
Let $\alpha\in[\rho,1]$ and $\delta:=1-\alpha$. For every prover and every session,
\[
\Pr[\text{the prover wins}]\le\eps_{\MCA}(\delta)+\eps_{\mathrm{SZ}}(\ell_\alpha)
+\eps_{\mathrm{OOD}}(\ell_\alpha)+(\alpha^-)^s .
\]
\end{theorem}

\begin{proof}
Condition on $z$; by \cref{prop:curves} (with base point $\zeta$, which is independent of
the session challenges), \cref{thm:session} applies to the session: except with
probability $\eps_{\MCA}(\delta)+\eps_{\mathrm{SZ}}(\ell_\alpha)+(\alpha^-)^s$, acceptance
implies $G(\tau)\ne\bot$, $G(\tau)\in\Lambda_\alpha(W_0)$ and $\gt=f_{G(\tau)}(\zeta+rT)$.
Then $G(\tau)(z)=f_{G(\tau)}(\zeta)=\gt(0)=y$. Outside the event of \cref{lem:ood},
$G(\tau)$ is therefore the unique list element with value $y$ at $z$, i.e.\
$G(\tau)=G_{z,y}$, and the prover does not win.
\end{proof}

\begin{corollary}[Anchored non-equivocation]\label{cor:anchored-eq}
For $m$ sessions on one commitment $(W_0,z,y)$, the probability that some session accepts
with $\gt$ different from $f_{G_{z,y}}(\zeta+rT)$, or with $G_{z,y}=\bot$, is at most
$\eps_{\mathrm{OOD}}(\ell_\alpha)+m\bigl(\eps_{\MCA}(\delta)+\eps_{\mathrm{SZ}}(\ell_\alpha)+(\alpha^-)^s\bigr)$.
In particular all accepted sessions open the same polynomial $G_{z,y}$, which is
determined by the commitment.
\end{corollary}

\begin{proof}
\Cref{lem:ood} is applied once; the remaining terms of \cref{thm:anchored} are applied to
each session.
\end{proof}

\begin{remark}[Anchored versus unique binding]\label{rem:anchored-vs-unique}
With $\alpha=\sqrt\rho+\eta$ and \cref{thm:mca-johnson}, \cref{thm:anchored} has query
factor $\sqrt\rho+\eta$ instead of $\frac{1+\rho}2$. The reference $G_{z,y}$ is fixed by the
commitment but chosen by the prover through $y$ among the list elements; in the
unique-decoding regime ($\alpha>\frac{1+\rho}2$) the list has at most one element and
$G_{z,y}\in\{F_\star,\bot\}$. The attack of \cref{thm:ub-tight} does not apply: its two
codewords $F$ and $F-H_Y$ differ at $z$ except with probability $\deg H_Y/|\K|$, so the check
$\gt(0)=y$ separates them.
\end{remark}

%=====================================================================
\section{Non-equivocation}
%=====================================================================

\begin{definition}[Equivocation game]\label{def:game-E}
The prover outputs $W_0:D_0\to\F$ once. Then two sessions $j=1,2$ of
\cref{prot:sfri} are run sequentially, both with initial oracle $W_0$, with fresh and
independent challenges $r^{(j)}$ and query points. The prover's strategy in session $2$
may depend on its entire view of session $1$, including the query points and the
verdict. Let $\tau_j$ be the commit transcript and $\Acc_j$ the acceptance event of
session $j$, and $G_j:=G(\tau_j)$ (\cref{def:opened}). The prover wins if
\[
\Acc_1\wedge\Acc_2\wedge\bigl(G_1=\bot\ \vee\ G_2=\bot\ \vee\ G_1\ne G_2\bigr).
\]
\end{definition}

\begin{theorem}[Non-equivocation]\label{thm:equivocation}
Let $\alpha\in[\rho,1]$, $\delta:=1-\alpha$, and
$M:=\max\bigl(\alpha^-,\ (\tfrac{1+\rho}2)^2\bigr)$. For every prover,
\[
\Pr[\text{the prover wins the equivocation game}]
\le2\bigl(\eps_{\MCA}(\delta)+\eps_{\mathrm{SZ}}(\ell_\alpha)\bigr)+M^{s}.
\]
\end{theorem}

\begin{proof}
For $j=1,2$, the prover restricted to session $j$ is a symbolic prover in the sense of
\cref{def:symbolic-game}: its strategy is randomized by its own coins and, for $j=2$, by
the independent coins of session $1$, while $r^{(j)}$ is fresh. Let $\mathsf E_j$ be the
event of \cref{thm:identification} for session $j$; then
$\Pr[\mathsf E_1\cup\mathsf E_2]\le2(\eps_{\MCA}(\delta)+\eps_{\mathrm{SZ}}(\ell_\alpha))$.

Write $\pi_j:=\pi^{\mathrm{sym}}_0(\tau_j)$ and let $\mathsf C$ be the event
$\neg\mathsf E_1\wedge\neg\mathsf E_2\wedge(G_1=\bot\vee G_2=\bot\vee G_1\ne G_2)$.

\emph{Claim: on $\mathsf C$, $\pi_1\pi_2\le M$.} If $\pi_1<\alpha$ then
$\pi_1\le\alpha^-$ and $\pi_1\pi_2\le\alpha^-$; likewise if $\pi_2<\alpha$. If
$\pi_1,\pi_2\ge\alpha$, then $\neg\mathsf E_j$ gives $G_j\ne\bot$ and
$\pi_j\le\agr(W_0,G_j)$, and $\mathsf C$ forces $G_1\ne G_2$. The sets
$A_j:=\{y\in D_1:W_0=G_j\text{ on }\fib(y)\}$ satisfy
$|A_1\cap A_2|<\rho n_1$ by \cref{lem:codeword-fibres}, so
$\pi_1+\pi_2\le(|A_1|+|A_2|)/n_1<1+\rho$, and
$\pi_1\pi_2\le(\frac{1+\rho}2)^2$ by the AM--GM inequality.

\emph{Conditioning.} Let $\mathcal F_1$ be the information up to the end of the commit
phase of session $1$, and $\mathcal F_2$ the information up to the end of the commit
phase of session $2$. Then $\Pr[\Acc_1\mid\mathcal F_1]=\pi_1^s$,
$\Pr[\Acc_2\mid\mathcal F_2]=\pi_2^s$, and $\mathsf C$, $\pi_1$ and $\pi_2$ are
$\mathcal F_2$-measurable. On $\mathsf C$ the claim gives
$\pi_2^s\le\min\bigl(1,(M/\pi_1)^s\bigr)$, a quantity $\mathcal F_1$-measurable
(set to $1$ if $\pi_1=0$). Hence
\begin{align*}
\Pr[\Acc_1\wedge\Acc_2\wedge\mathsf C]
&=\mathbb E\bigl[\mathbf 1_{\Acc_1}\mathbf 1_{\mathsf C}\,\pi_2^s\bigr]
\le\mathbb E\bigl[\mathbf 1_{\Acc_1}\min(1,(M/\pi_1)^s)\bigr]\\
&=\mathbb E\bigl[\pi_1^s\min(1,(M/\pi_1)^s)\bigr]\le M^s ,
\end{align*}
where the middle equality conditions on $\mathcal F_1$ (the event $\Acc_1$ depends only on
the query points of session $1$, which are independent of $\mathcal F_1$). Since the
winning event is contained in $\mathsf E_1\cup\mathsf E_2\cup(\Acc_1\wedge\Acc_2\wedge\mathsf C)$,
the theorem follows.
\end{proof}

\begin{remark}\label{rem:equivocation-regimes}
In the unique-decoding regime ($\alpha$ just above $\frac{1+\rho}2$) one has
$M=\alpha^-\le\frac{1+\rho}2$ and \cref{thm:equivocation} gives nothing beyond
\cref{thm:unique-binding}. Up to the Johnson radius ($\alpha=\sqrt\rho+\eta$),
$M=\max(\sqrt\rho+\eta,(\frac{1+\rho}2)^2)$, which is strictly smaller than
$\frac{1+\rho}2$. For $\rho\in\{\frac12,\frac14,\frac18\}$ one has
$\sqrt\rho>(\frac{1+\rho}2)^2$, so $M=\sqrt\rho+\eta$; for $\rho=\frac1{16}$,
$(\frac{17}{32})^2\approx0.2822>\frac14+\eta$ whenever $\eta<0.0322$, so
$M=(\frac{17}{32})^2$.
\end{remark}

%=====================================================================
\section{Summary of query factors}
%=====================================================================

\begin{center}
\renewcommand{\arraystretch}{1.3}
\begin{tabular}{lll}
\hline
Notion & Query factor (per path) & Coding input\\
\hline
Unique binding (\ref{thm:unique-binding}) & $\frac{1+\rho}2$ & MCA, unique decoding (proved)\\
Proximity (\ref{thm:proximity}) & $\sqrt\rho+\eta$ & MCA up to Johnson\\
Promise binding (\ref{thm:promise}) & $\max(\sqrt\rho+\eta,\ \delta_0+\rho)$ & MCA up to Johnson\\
Session uniqueness (\ref{thm:session}) & $\sqrt\rho+\eta$ & MCA up to Johnson\\
Anchored binding (\ref{thm:anchored}) & $\sqrt\rho+\eta$ & MCA up to Johnson\\
Non-equivocation (\ref{thm:equivocation}) & $\max(\sqrt\rho+\eta,\ (\frac{1+\rho}2)^2)$ & MCA up to Johnson\\
\hline
\end{tabular}
\end{center}
Each bound adds $\eps_{\MCA}$ and a root-count term of order $\ell_\alpha^2d/|\K|$.
